\begin{resumo}
A ISO 20816-3 avalia a vibração de máquinas industriais pela velocidade eficaz medida na carcaça e pela mudança desse valor, e um acelerômetro de sistemas microeletromecânicos (MEMS) de consumo custa de uma a duas ordens de grandeza menos que o acelerômetro piezoelétrico usado como referência nessa medição. Em dois nós de baixo custo publicados que usam o detector Isolation Forest, o falso alarme foi de \SI{22}{\percent} a \SI{38}{\percent} da operação normal, e nenhum dos trabalhos consultados avaliou o falso alarme entre rodadas de medição da mesma máquina numa instalação industrial. Este trabalho avaliou se um sistema de aquisição de baixo custo, que transmite o sinal bruto de vibração, permite detectar, com um modelo treinado só com a operação normal de cada máquina, mudanças de \SI{0,7}{\milli\meter\per\second}, \SI{25}{\percent} do limite superior da zona de vibração aceitável da norma, com falso alarme limitado entre rodadas. Um transmissor com ESP32-C3 e módulo MPU-6050 adquiriu janelas de \SI{0,512}{\second} a \SI{1}{\kilo\hertz} e as enviou por rádio LoRa a um \textit{gateway}, que as repassou a uma interface de programação de aplicações, e o sistema foi verificado por testes automatizados e por ensaios de ponta a ponta. Com impostos, os componentes do transmissor e do \textit{gateway} custam R\$~\num{176,04}, sem a bateria. Em 2 de novembro de 2026, ele adquiriu \campoJanelasOperacao\ janelas em \campoEnsaios\ ensaios de \campoMaquinas\ máquinas de uma instalação industrial, todas em operação normal. Quatro hipóteses foram avaliadas com critérios registrados antes do cálculo e depois de análises exploratórias, a quarta por emenda: falso alarme entre duas rodadas da mesma máquina (H1), menor mudança detectável na velocidade eficaz abaixo de \SI{0,7}{\milli\meter\per\second} (H2), detecção de uma assinatura de desbalanceamento somada digitalmente a janelas reais (H3) e detecção de desbalanceamento real numa base pública de laboratório (H4). As janelas analisadas chegaram completas, e nenhuma hipótese foi corroborada. H1 ficou \hUmVeredito, com \hUmCalandraIFDireto~\% das janelas da segunda rodada da calandra marcadas, e H2 foi \hDoisVeredito, com menor mudança detectável de até \hDoisMmdCalandraDois~mm/s. H3 foi \hTresVeredito, com no máximo \hTresFracaoMaxCalandra~\% das janelas marcadas até \SI{2,8}{\milli\meter\per\second} de assinatura, na calandra, que já tinha \apUmFracaoSemAssinaturaCalandra~\% marcadas sem assinatura. H4 foi \hQuatroVeredito, com \hQuatroFracaoZero~\% de falso alarme depois da remontagem do suporte do desbalanceamento, embora de \hQuatroFracaoTres~\% a \hQuatroFracaoQuatro~\% das janelas com desbalanceamento tenham sido marcadas. Em análises posteriores, H1 seria corroborada com \apTresVereditoUmCorroborada\ das \apTresSementes\ sementes do gerador de números aleatórios do detector, a menor mudança detectável ficou de \apCincoRazaoMin\ a \apCincoRazaoMax\ vezes a do sensor em repouso, e o escore do Isolation Forest ficou abaixo do limiar em três das quatro máquinas mesmo com os descritores alterados pela assinatura fora da faixa do treino. Nesses dados, a detecção foi limitada pelo escore do detector, pela linha de base de uma rodada, com \campoJanelasGuilhotinaUm\ a \campoJanelasCalandraUm\ janelas por máquina, e pela variação entre janelas medida nas máquinas, que reúne a vibração da máquina e a resposta da fixação do sensor.

\vspace{\onelineskip}
\noindent
\textbf{Palavras-chave}: detecção de falhas; manutenção preditiva; análise de vibração; acelerômetro MEMS; Isolation Forest.
\end{resumo}
